\begin{table}[H]
\centering
\caption{Contribución de cada actividad a la PTF total por subperiodo}
\label{tab:contribucion_ptf_subperiodos}
\footnotesize
\setlength{\tabcolsep}{3pt}
\begin{tabular}{p{5.3cm}*{5}{>{\raggedleft\arraybackslash}p{1.45cm}}}
\toprule
Actividad & 2005--2010 & 2011--2015 & 2016--2019 & 2020--2024 & 2005--2024 \\
\midrule
Agropecuario & -0,07 & 0,12 & 0,09 & 0,11 & 0,06 \\
Minería & 0,01 & -0,19 & -0,15 & -0,15 & -0,11 \\
Manufactura & -0,02 & -0,16 & -0,02 & -0,03 & -0,06 \\
Servicios públicos & -0,01 & -0,04 & -0,02 & 0,03 & -0,01 \\
Construcción & 0,06 & -0,02 & -0,08 & -0,20 & -0,05 \\
Comercio y hotelería & 0,09 & 0,22 & 0,08 & 0,14 & 0,13 \\
Transporte y comunicaciones & -0,04 & 0,08 & 0,04 & 0,00 & 0,02 \\
Financieras & -0,22 & -0,04 & -0,17 & -0,17 & -0,15 \\
Servicios sociales y personales & -0,03 & 0,05 & 0,19 & 0,28 & 0,11 \\
\midrule
\textbf{Total de la economía} & \textbf{-0,23} & \textbf{0,02} & \textbf{-0,04} & \textbf{0,00} & \textbf{-0,07} \\
\bottomrule
\end{tabular}
\vspace{0.15cm}
\begin{minipage}{0.96\textwidth}
\footnotesize \textit{Nota:} promedio anual en puntos porcentuales. Cada celda es el promedio del producto entre el peso anual de la actividad y la tasa anual de crecimiento de su PTF. Los subperiodos tienen distinta duración.
\par\textit{Fuente:} cálculos del CJC con base en DANE, anexo PTF 2025.
\end{minipage}
\end{table}
